\documentclass[10pt]{cv}

% \usepackage[hmargin=0.5cm, vmargin=0.125cm]{geometry} % Document margins
\usepackage{fontawesome} % Required for symbols in the colored box
\usepackage{amssymb} % \mathbb{X} doubles as the X (formerly Twitter) logo
% \usepackage[usenames,dvipsnames]{xcolor} % Allows the definition of hex colors

\usepackage[utf8]{inputenc}
\usepackage[T1]{fontenc}

% Colors for links, text and headings
\usepackage{hyperref}

\hypersetup{colorlinks,breaklinks, urlcolor=LinkCyan, linkcolor=LinkCyan} % Set up links and colors

% \usepackage{fancyhdr}
% \pagestyle{fancy}
% \fancyhf{}


% \renewcommand{\headrulewidth}{0pt} % Get rid of the default rule in the header

\usepackage{titlesec} % Allows creating custom \section's
\usepackage{adforn} % Ornaments

\usepackage[export]{adjustbox}
\usepackage{pdfcomment}

\titlespacing{\section}{0pt}{0pt}{5pt} % Spacing around titles

%%%%%%%%%%%%%%%%%%%%%%%%%%%%%%%%%%%%%%%%%
%%                                     %%
%%   SET PER COMPANY BEFORE SENDING!   %%
%%                                     %%
%%%%%%%%%%%%%%%%%%%%%%%%%%%%%%%%%%%%%%%%%
%%                                     %%
%%  This variable is used in the       %%
%%  utm_content parameter. Every       %%
%%  company that gets the CV should    %%
%%  have a different value. This way   %%
%%  we are able to track which         %%
%%  interviewers from which companies  %%
%%  actually read the CV, when,        %%
%%  and what blog posts they read.     %%
%%                                     %%
%%  Don't edit it here - build with:   %%
%%    ./build.sh <company> [...]       %%
%%                                     %%
%%%%%%%%%%%%%%%%%%%%%%%%%%%%%%%%%%%%%%%%%
\providecommand{\uniquecampaign}{unmodified}

% {post path}{utm_medium}{text}
% e.g: \bloglink{2019/03/trick-or-trait.html}{extras-sidebar}{text to displayed (passed to href as is)}
\newcommand{\bloglink}[3][]{
    \href{https://blog.hochgi.com/#1?utm_source=cv&utm_medium=#2&utm_campaign=bonus_points_for_checking_my_blog&utm_content=\uniquecampaign}{#3}
}

% Key term (tech, skill or domain) highlighted in the text
\newcommand{\kw}[1]{\textbf{#1}}

\begin{document}

    %%%%%%% Printable mode
    \PrintableMode

    %----------------------------------------------------------------------------------------
    %	TOP BAR
    %----------------------------------------------------------------------------------------
    \begin{TopBar}{\ColorTextSide}
        \Name{\ColorHighlight}{\textit{Gilad Hoch}}{\small\color{MainOrange}Backend \& AI Tech Lead}{}

        \SimpleSeparator{\ColorHighlight}
        \RaggedRight\small{
            \textbf{About me:} Backend engineer and tech lead with 14+ years of experience in \kw{distributed systems}, \kw{big data} and \kw{API design},
            now leading AI engineering: agents, harnesses, and making AI-assisted development reliable \emph{and} affordable.
            Thorough, opinionated about how software should be built, and always happy to learn something new.
        }

        \SimpleSeparator{\ColorHighlight}
        \begin{TripleColumns}
            \begin{ItemList}{\ColorHighlight}
                \item [\Large\faEnvelopeO]      \href{mailto:giladhoch@gmail.com}{giladhoch@gmail.com} % Email address
                \item [\Large\faPhone]          +972 58 6204060 % Phone
                \item [\Large\faLinkedinSquare] \href{https://www.linkedin.com/in/gilad-hoch-6632a138}{LinkedIn profile} % Linkedin
            \end{ItemList}
            \nextcolumn
            \begin{ItemList}{\ColorHighlight}
                \item [\Large\faMapMarker] Rehovot, Israel                                       % Location
                \item [\Large\faLaptop]    \bloglink{top-header}{blog.hochgi.com}                 % Website
                \item [\Large\faGithub]    \href{https://github.com/hochgi}{github.com/hochgi}    % Github
            \end{ItemList}
            \nextcolumn
            \begin{ItemList}{\ColorHighlight}
                \item [\Large$\mathbb{X}$]     \href{https://x.com/hochgi}{X profile}                                 % X (formerly Twitter)
                \item [\Large\faStackExchange] \href{https://stackexchange.com/users/482262/gilad-hoch?tab=accounts}{Stack Exchange profile} % Stackexchange
                \item [\Large\faStackOverflow] \href{https://stackoverflow.com/users/896557/gilad-hoch}{Stack Overflow profile} % Stackoverflow
            \end{ItemList}
        \end{TripleColumns}
    \end{TopBar}


    %----------------------------------------------------------------------------------------
    %	EXPERIENCE
    %----------------------------------------------------------------------------------------
    \MainTitle{\ColorHighlight}{\ColorTextMain}{Experience}

    \begin{flushleft}


        %%%%%%%%%%%%%
        % Versatile %
        %%%%%%%%%%%%%
        \Company{Versatile}{Feb 2025 -- Present}

        \Role{AI Tech Lead}{Aug 2026 -- Present}
        Leading the company's AI initiatives: agentic coding infrastructure, AI tooling and \kw{AI cost optimization}.
        Started in early 2026 as a side initiative alongside my backend work; a formal role since Aug 2026.
        \begin{itemize}
            \item Built \emph{Van Damme}, a Slack-summoned general-purpose \kw{AI agent} that authored 250+ PRs and reviewed 330+
                  across 30+ repositories within months, becoming part of everyday work across the company. It debugs, analyzes data,
                  and opens \& reviews PRs, through specialized agents with their own skills, tools and integrations
                  (GitHub, Jira, Databricks, HubSpot, Notion, Coralogix, the company ERP and more).
            \item Cut monthly AI spend almost 10$\times$ on a reference seat while doing more than double the work: a central \kw{LiteLLM}
                  gateway with per-user keys \& cost reporting routes the grind to cheaper models (xAI Grok, valarhq.ai, Bedrock)
                  and keeps frontier models for judgment, backed by a seat policy across tools.
            \item Introduced \kw{harness engineering}: a five-phase spec-to-ship \kw{agent pipeline} (spec, red tests, code to green, review, verify)
                  with per-phase permissions, so no agent grades its own homework. Rolled it out to key repositories, coached teams to adopt it,
                  and kept it in sync across \kw{Claude Code}, \kw{Cursor} \& \kw{OpenCode} via generated configs and a drift-checking pre-push hook.
            \item Evangelized AI adoption through Slack channels, tech talks and mentoring.
        \end{itemize}

        \Role{Sr. Backend Engineer}{Feb 2025 -- Present}
        Joined the backend team (\kw{node.js} / \kw{TypeScript} microservices), took ownership of the reporting \& widgets domain,
        and continue hands-on backend work alongside the AI role.
        \begin{itemize}
            \item Took PagerDuty alerts from almost daily to none and retired Power BI (and its licensing costs) by designing and leading the sunset
                  of the legacy reports pipeline, a screenshot-based flow wired into many services. Replaced it with native on-demand widgets,
                  now live in the web app and mobile clients.
            \item Created \href{https://github.com/hochgi/test-kit}{\emph{test-kit}} (open-sourced), a \kw{component-testing} framework for \kw{TypeScript}
                  services adopted across the backend: test doubles with probes that fully control them, backed by in-memory \kw{Postgres}, \kw{MySQL},
                  \kw{Redis}, \kw{S3}, \kw{SQS} \& \kw{Kafka}. Also wrote the \texttt{/component-testing} agent skill that makes proper component
                  testing the easy path for humans and agents alike.
        \end{itemize}



        %%%%%%%%%%%
        % Hunters %
        %%%%%%%%%%%
        \Company{Hunters}{Aug 2023 -- Feb 2025}

        \Role{Tech Lead}{}
        Mentored the team and provided technical guidance, focusing on ingestion \& detection. Played a key role in stabilizing a complex
        health monitoring system, resolving major \kw{concurrency} issues and cutting metrics reporting time from \emph{minutes} to single-digit
        \emph{milliseconds}; built a live visualization of alerting metrics that significantly shortened root-cause investigations; introduced
        \kw{OpenAPI} documentation with auto-generated clients; and led a \kw{Spark} upgrade with an overhaul of the build \& packaging process.
        Became the technical focal point for multiple \kw{Scala} \& \kw{ZIO} services, simplifying and documenting them so the team could change them with confidence.



            %%%%%%%%%%%%%%%
            % SparkBeyond %
            %%%%%%%%%%%%%%%
            \Company{SparkBeyond}{Jan 2019 -- Jul 2023}

            \Role{Team Lead}{Dec 2022 -- Jul 2023}
            Modernized the legacy product's API, rewriting it with \href{https://tapir.softwaremill.com}{tAPIr} to expose \kw{OpenAPI} docs and enable
            integrations. Then led the team into the product's core engine for the first time (segmentation analysis): pitched the approach to product
            \& VPs and ran a swift, productive onboarding.

            \Role{Tech Lead}{Apr 2021 -- Dec 2022}
            Recruited and onboarded a new team for a greenfield project continuing my IC work, and led it technically: caching proxies for heavy
            NLP models, a throttler \& dispatcher (built on \kw{akka-streams}) for bottleneck services such as \kw{Elasticsearch}, a data warehouse
            on \kw{Spark}, \kw{Cassandra}, \kw{Elasticsearch} \& \kw{Kafka}, and a business-logic service on \kw{Postgres}.

            \Role{IC}{Jan 2021 -- Apr 2021}
            Proposed the architecture and work plan to commercialize a labs POC, and laid the product's foundations together with a DevOps engineer.

            \Role{Sr. Backend Engineer}{Jan 2019 -- Jan 2021}
            On the backend platform team: full \kw{Parquet} support as an alternative to CSV, a modernized telemetry pipeline, a simplified build,
            a cross-version \kw{Scala} SDK for satellite projects, and most of the PII \& bleaching logic. One of the main reviewers and mentors;
            advised the labs team on architecture \& design.

            \Role{Additional roles}{2020 -- 2023}
            Backend Guild Manager \& Guilds Manager: meetups, tech talks, mentoring sessions, and a full-week company-wide \kw{Scala} course
            delivered by \href{https://ziverge.com/}{Ziverge}. Intern Lead (2022): took an intern with no prior knowledge to all major goals.

            %%%%%%%%%%%%%%%%%%%%%%%%%%%%
            % Aston Capital Management %
            %%%%%%%%%%%%%%%%%%%%%%%%%%%%
            \Company{Aston Capital Management}{Apr 2018 -- Dec 2018}

            \Role{Software Architect}{}
            Forex trading system: \kw{DB schemas}, specs \& formats; an \kw{akka-streams} pipeline exporting large volumes of data to \kw{AWS S3};
            \kw{REST} APIs for data access; ETL for the data science team; and live data \kw{streaming} \& monitoring.

            %%%%%%%%%%%%%%%%%%%
            % Thomson Reuters %
            %%%%%%%%%%%%%%%%%%%
            \Company{Thomson Reuters}{Jun 2012 -- Apr 2018}

            \Role{Software Architect}{Jan 2018 -- Apr 2018}
            In-house consulting on architecture, design and \kw{data modeling} for teams building on CM-Well.

            \Role{Sr. Software Developer}{Jul 2014 -- Jan 2018}
            Matured CM-Well (built with \kw{Scala}, \kw{akka} \& \kw{Play}) into a standard Thomson Reuters service and an open-source project on GitHub.

            \Role{Software Developer (contractor)}{Jul 2012 -- Jul 2014}
            Web service APIs, backend refactoring to a new architecture, and migrating the build from Maven to \kw{sbt}.
    \end{flushleft}

    %----------------------------------------------------------------------------------------
    %	EDUCATION
    %----------------------------------------------------------------------------------------
    %%% Education
    \MainTitle{\ColorHighlight}{\ColorTextMain}{Education}
    \Experience
        {\ColorHighlight}
        {Graduate studies in Computer Science (M.Sc. track)}
        {Bar-Ilan University}
        {Sep 2012 -- Jan 2015}
        {
            Graduate coursework toward an M.Sc. (degree not completed). Area of study: Algorithms, Data structures \& Complexity.
        }
    \Experience
        {\ColorHighlight}
        {B.Sc. in Computer Science}
        {Bar-Ilan University}
        {Sep 2009 -- Jun 2012}
        {
            Was accepted into the Microsoft Excellence Program and received a scholarship. Participated in the RoboCup project: the Bar-Ilan University RoboCup Soccer Team (BURST) was founded in 2009 and was the first team from Israel to compete in the RoboCup senior league. The team consisted of a few selected undergraduate students \& researchers.
        }
    \Experience
        {\ColorHighlight}
        {Microsoft Excellence Program}
        {Microsoft R\&D Center, Herzliya}
        {Sep 2009 -- Jun 2012}
        {
            Participated in the program led by Microsoft and attended enrichment courses on various topics, e.g.: \\
            \scriptsize{.NET environment \& C\#, XNA, Kinect SDK, Mobile (Windows Phone) Development, Azure \& Big Data}
        }


    \begin{flushleft}
        \normalsize{\color{MidGrey}\textbf{3D Animation course\hfill{\small\color{\ColorHighlight}Apr 2006 -- May 2007}}}\\\vspace{-.4ex}
        \ssmall{\color{\ColorHighlight}\hrulefill}
        \small{\color{MidGrey}\textbf{Adventure Publishing Corp.}}\\
        \small{\color{MidGrey}A year-long course in 3D \& animation studies (modeling, texturing, rigging, animating).} \\
    \end{flushleft}


    %----------------------------------------------------------------------------------------
    %	FOOTER BAR (bottom of the last page)
    %----------------------------------------------------------------------------------------
    \begin{FooterBar}
        \begin{TripleColumns}
            \raggedright
            \TextSeparatorBis{\ColorHighlight}{\faCaretRight \hspace{1pt} Talks}{}
            \faYoutubePlay~\href{https://youtu.be/P21eSH2-0HY}{Concurrent Scala} \\
            \faYoutubePlay~\href{https://youtu.be/pB7EMXtfQUY}{Parsing - beyond regex} \\
            \faDesktop~\href{https://hochgi.github.io/how-to-ai-in-versatile/}{How to AI in Versatile} {\small\color{MidGrey}(slides)} \\
            \faDesktop~\href{https://docs.google.com/presentation/d/1lweerpmM9owiQNuImpmT9ROU4zRlym_jIdzbVkLJ-JQ/edit?usp=sharing}{Component Testing} {\small\color{MidGrey}(slides)} \\[0.5ex]
            {\footnotesize\color{MidGrey}Most of my talks weren't recorded or published openly.}
            \nextcolumn
            \raggedright
            \TextSeparatorBis{\ColorHighlight}{\faCaretRight \hspace{1pt} Open Source}{}
            \faGithub~\href{https://github.com/hochgi?tab=repositories&sort=stargazers}{My repositories} \\
            \faGithub~\href{https://github.com/hochgi/test-kit}{test-kit} \\
            \faGithub~Contributions to \href{https://github.com/pulls?utf8=\%E2\%9C\%93&q=is\%3Apr+author\%3Ahochgi+org\%3Aakka}{Akka} \& \href{https://github.com/pulls?utf8=\%E2\%9C\%93&q=is\%3Apr+author\%3Ahochgi+org\%3Azio}{ZIO} \\
            \faGithubSquare~\href{https://api.github.com/search/issues?q=author\%3Ahochgi+type\%3Apr&per_page=100&page=1}{All my PRs} \\[0.5ex]
            {\footnotesize\color{MidGrey}A GitHub API link, 100 results per page; a little
                \pdftooltip{\texttt{\textcolor{MainOrange}{jq}}~\&~\texttt{\textcolor{MainOrange}{bash}}~magic}{
                    cat <(jq '.items\\[\\]' page1.json) <(jq '.items\\[\\]' page2.json) <(jq '.items\\[\\]' page3.json) | jq -sr 'map(.repository\_url)\\[\\]' | awk -F/ '\{print \$5\}' | sort | uniq
                } over all pages lists the orgs: \emph{akka, CM-Well, coursier, OpenAPITools, playframework, sbt, scala, zio, \dots}}
            \nextcolumn
            \raggedright
            \TextSeparatorBis{\ColorHighlight}{\faCaretRight \hspace{1pt} Writing}{}
            \faRss~\bloglink{extras-footer}{blog.hochgi.com} \\[0.5ex]
            {\footnotesize\color{MidGrey}On software engineering, Scala \& functional programming, and engineering culture.}
        \end{TripleColumns}
    \end{FooterBar}

\end{document}
